\section{Gestión del proyecto y trazabilidad}
\label{sec:gestion}

El trabajo se organizó en 32 \emph{Issues} con criterios de aceptación
explícitos, 29 de ellos cerrados, y 31 \emph{Pull Requests} integrados
---todos con al menos una revisión del otro integrante y 11 con dos--- sobre
149 commits. Dos \emph{Issues} de diseño precedieron a la implementación:
uno fijó el formato del segmento y el protocolo de inserción y extracción, y
el otro el registro de participantes, el protocolo de cierre y el formato de
los registros. Ninguno de los dos se implementó antes de estar revisado. Las
decisiones se anotaron en la bitácora a medida que se tomaban, cada entrada
en el \emph{Pull Request} que la implementó~\cite{Bitacora2026}, no al final.

\subsection{Responsabilidades}

El reparto siguió los módulos, no las líneas.

\textbf{Emmanuel Barrantes} trabajó el formato del segmento y el protocolo
de inserción y extracción, el buffer circular, la capa de sincronización, el
reloj monotónico con el generador pseudoaleatorio y la pausa exponencial,
las métricas con el histograma de latencia, el payload verificable, el
verificador \code{verify}, el ejecutable \code{producer} y la primera mitad
del \code{finisher}, el resumen final, el script del experimento y su
análisis, las pruebas de conservación, escala, espera bloqueante y cierre, y
el objetivo de los sanitizers.

\textbf{Isaac Moreno} trabajó el repositorio con su Makefile y su README, el
registro de participantes junto con el protocolo de cierre y el formato de
los registros, el módulo de memoria compartida con el \code{creator}, el
análisis y validación de argumentos, el registrador CSV, el ejecutable
\code{consumer}, la segunda mitad del \code{finisher} con la espera de
retiradas y la liberación, las pruebas de validación, aislamiento y carreras
de ciclo de vida, y la verificación final del código contra el enunciado.

En este informe, las secciones~\ref{sec:arquitectura} a~\ref{sec:pruebas}
son de Isaac y las~\ref{sec:resultados} a~\ref{sec:decisiones} de Emmanuel;
el resto es conjunto.

\subsection{Issues, ramas y revisiones}

Cada cambio entró por una rama con nombre derivado de su \emph{Issue} y un
\emph{Pull Request} que declara la validación ejecutada y queda a revisión
del otro integrante. Los commits pertenecen a quien hizo el trabajo; no se
reescribió historia compartida, no se sustituyó el repositorio al final y
ningún cambio entró como un único commit.

Las revisiones no fueron de forma. Una encontró que la fila
\code{lock\_acquire} del resumen declaraba incluir la readquisición del
mutex al despertar de la espera, cuando el tiempo de bloqueo se mide una
sola vez alrededor de la toma inicial y esa readquisición queda
contabilizada en el tiempo de espera. Los datos lo confirmaron ---la nota
implicaba unos 60~\textmu s y lo medido son 0.89 a 1.74~\textmu s--- y se
corrigió la nota junto con la prueba que la aseveraba. Lo detectó el cruce
de la documentación con el código, no una prueba, que habría pasado igual.

Dos incidentes del flujo quedaron en el historial en vez de borrarse: un
\emph{Pull Request} integrado contra la rama equivocada, que dejó su trabajo
fuera de la rama principal y hubo que rehacerlo contra ella, y una colisión
de numeración en la bitácora cuando dos ramas reclamaron la misma entrada a
la vez.
